\subsection{Метод квадратичной аппроксимации}
\label{subsec:quad-approx}

\begin{idea}
В тетради название метода пропущено. По содержанию это \textbf{метод
квадратичной аппроксимации} — прямой аналог метода Ньютона для задач с
ограничениями: там, где метод условного градиента (\S\ref{subsec:cond-gradient})
минимизирует линейную модель, здесь минимизируется квадратичная. На нём же
основано семейство методов последовательного квадратичного программирования.
\end{idea}

\[
  \begin{cases}
    \varphi(x) \to \min , & x \in S , \\
    \varphi(x) \ \text{— выпуклая, дважды дифференцируемая} , \\
    S \ \text{— выпуклый компакт} .
  \end{cases}
\]

\begin{algo}[title={Метод квадратичной аппроксимации}]
\textbf{I. Начальный этап.} Задать $x_0 \in S$ — начальное приближение,
$k = 0$; итерации вида $x_{k+1} = x_k + \alpha_k s_k$, $\alpha_k > 0$.

\medskip
\textbf{II. $k$-й шаг.} Рассмотрим квадратичную аппроксимацию в окрестности
$x_k$:
\begin{steps}
  \item Положим
        \[
          \psi_k(x) = \nabla\T\varphi(x_k)\,(x - x_k)
            + \tfrac{1}{2}\,(x - x_k)\T H(x_k)\,(x - x_k)
        \]
        и решим вспомогательную задачу
        \[
          \psi_k \to \min , \qquad x \in S .
        \]
        Пусть $y_k$ — её решение; так как $x_k \in S$ и $\psi_k(x_k) = 0$,
        \[
          \eta_k = \psi_k(y_k) \le \psi_k(x_k) = 0 .
        \]
        В качестве направления берём $s_k = y_k - x_k$.
  \item Шаг выбирается как в методе второго порядка для задач без ограничений:
        заданы $\alpha_0 = 1$, $\eps \in (0,1)$, $\lambda \in (0,1)$, и
        $\alpha_k = \alpha_0 \lambda^{\,i_k}$ — первое значение, удовлетворяющее
        условию дробления
        \begin{equation}\label{eq:qa-armijo}
          \varphi(x_k + \alpha s_k) - \varphi(x_k) \ \le\ \alpha\,\eps\,\eta_k
        \end{equation}
        (либо $\varphi(x_{k+1}) \to \min_{\alpha > 0}$ — наискорейший спуск
        вдоль $s_k$).
\end{steps}
\end{algo}

\subsubsection*{Вспомогательные факты}

\begin{lemma}\label{lem:qa-basic}
В условиях задачи:
\begin{enumerate}[label=$\arabic*^\circ$,left=0pt]
  \item $x_k + \alpha s_k \in S$ при всех $\alpha \in [0,1]$;
  \item $\nabla\T\varphi(x_k)\, s_k \le \eta_k \le 0$;
  \item если дополнительно $\varphi$ сильно выпукла с константой $m$, то
        \[
          \eta_k \ \le\ -\tfrac{m}{2}\,\norm{s_k}^2 .
        \]
\end{enumerate}
\end{lemma}

\begin{proof}
$1^\circ$ $x_k + \alpha s_k = (1-\alpha)x_k + \alpha y_k$ — выпуклая комбинация
точек $x_k, y_k \in S$.

$2^\circ$ По определению $\eta_k = \psi_k(y_k)
= \nabla\T\varphi(x_k)s_k + \tfrac12 s_k\T H(x_k) s_k$, а второе слагаемое
неотрицательно в силу выпуклости $\varphi$ ($H \succeq 0$). Неравенство
$\eta_k \le 0$ — из $\psi_k(x_k) = 0$.

$3^\circ$ Функция $\psi_k$ сильно выпукла с той же константой $m$, а $y_k$ —
её минимум на $S$, поэтому $\nabla\T\psi_k(y_k)(x_k - y_k) \ge 0$. Тогда
\[
  \underbrace{\psi_k(x_k)}_{=\ 0}
  \ \ge\ \psi_k(y_k) + \underbrace{\nabla\T\psi_k(y_k)(x_k - y_k)}_{\ge\ 0}
    + \tfrac{m}{2}\norm{x_k - y_k}^2
  \ \ge\ \eta_k + \tfrac{m}{2}\norm{s_k}^2 . \qedhere
\]
\end{proof}

\begin{lemma}\label{lem:qa-step}
Дробление~\eqref{eq:qa-armijo} конечно, и принятый шаг отделён от нуля:
существует $\bar\alpha > 0$ с $\alpha_k \ge \min\{1, \bar\alpha\,\abs{\eta_k}\}$.
\end{lemma}

\begin{proof}
Гессиан непрерывен на компакте $S$, значит ограничен: $\norm{H(x)} \le M$.
Тогда
\[
  \varphi(x_k + \alpha s_k) \ \le\ \varphi(x_k)
    + \alpha \nabla\T\varphi(x_k) s_k + \tfrac{M}{2}\alpha^2 \norm{s_k}^2
  \ \overset{\text{лемма }\ref{lem:qa-basic}}{\le}\
  \varphi(x_k) + \alpha\eta_k + \tfrac{M}{2}\alpha^2\norm{s_k}^2 .
\]
Для выполнения~\eqref{eq:qa-armijo} достаточно
\[
  \alpha\eta_k + \tfrac{M}{2}\alpha^2\norm{s_k}^2 \le \alpha\eps\eta_k
  \qquad\Longleftrightarrow\qquad
  \alpha \ \le\ -\frac{2(1-\eps)\,\eta_k}{M \norm{s_k}^2} ,
\]
и правая часть положительна при $\eta_k < 0$. Значит после конечного числа
дроблений условие выполнится. Если же принятый шаг не равен $\alpha_0$, то
предыдущее пробное значение $\alpha_k/\lambda$ было отвергнуто, откуда
\[
  \alpha_k \ >\ -\frac{2\lambda(1-\eps)\,\eta_k}{M \norm{s_k}^2}
  \ \ge\ \frac{2\lambda(1-\eps)}{M C^2}\, \abs{\eta_k} ,
\]
где $C = \operatorname{diam} S \ge \norm{s_k}$.
\end{proof}

\subsubsection*{Теоремы о сходимости}

\begin{theorem}\label{thm:qa-1}
Пусть $\varphi(x)$ удовлетворяет условиям постановки задачи и $\{x_k\}$
построена по описанной схеме. Тогда независимо от выбора начального
приближения последовательность $\varphi(x_k)$ монотонно убывает и
\[
  \lim_{k \to \infty} \varphi(x_k) = \varphi(x_*) = \min_{x \in S}\varphi(x) .
\]
\end{theorem}

\begin{idea}[title={Доказательство (на лекции не приводилось)}]
\begin{proof}
\textbf{Шаг 1: убывание и $\eta_k \to 0$.} Из~\eqref{eq:qa-armijo}
\[
  \varphi(x_k) - \varphi(x_{k+1}) \ \ge\ \eps\,\alpha_k\,\abs{\eta_k} \ \ge\ 0 ,
\]
то есть $\varphi(x_k)$ не возрастает; она ограничена снизу ($S$ — компакт,
$\varphi$ непрерывна), значит сходится, и $\varphi(x_k) - \varphi(x_{k+1}) \to 0$.
По лемме~\ref{lem:qa-step}
\[
  \varphi(x_k) - \varphi(x_{k+1})
  \ \ge\ \eps \min\Big\{ \abs{\eta_k},\ \tfrac{2\lambda(1-\eps)}{MC^2}\eta_k^2 \Big\} ,
\]
и правая часть стремится к нулю только если $\eta_k \to 0$.

\medskip
\textbf{Шаг 2: предельная точка оптимальна.} Пусть $x_{k_j} \to \bar{x}$
(существует по компактности $S$). По определению $y_k$ как минимума $\psi_k$
на $S$,
\[
  \eta_k \ \le\ \psi_k(x)
  = \nabla\T\varphi(x_k)(x - x_k) + \tfrac12 (x - x_k)\T H(x_k)(x - x_k)
  \qquad \forall x \in S .
\]
Переходя к пределу по $k_j$ и пользуясь $\eta_{k_j} \to 0$ и непрерывностью
$\nabla\varphi$, $H$:
\[
  0 \ \le\ \nabla\T\varphi(\bar{x})(x - \bar{x})
    + \tfrac12 (x - \bar{x})\T H(\bar{x})(x - \bar{x})
  \qquad \forall x \in S .
\]
Подставим $x = \bar{x} + t(z - \bar{x})$, где $z \in S$, $t \in (0,1]$ (такая
точка лежит в $S$ по выпуклости):
\[
  0 \ \le\ t\,\nabla\T\varphi(\bar{x})(z - \bar{x})
    + \tfrac{t^2}{2}(z - \bar{x})\T H(\bar{x})(z - \bar{x}) .
\]
Разделив на $t > 0$ и устремив $t \to 0$, получаем
\[
  \nabla\T\varphi(\bar{x})\,(z - \bar{x}) \ \ge\ 0 \qquad \forall z \in S ,
\]
а это — необходимое и (в силу выпуклости $\varphi$) достаточное условие
минимума. Значит $\bar{x}$ — точка минимума, и
$\varphi(\bar{x}) = \varphi(x_*)$.

\medskip
\textbf{Шаг 3.} Последовательность $\varphi(x_k)$ монотонна и имеет предел;
по шагу~2 вдоль подпоследовательности этот предел равен $\varphi(x_*)$,
следовательно $\varphi(x_k) \to \varphi(x_*)$.
\end{proof}
\end{idea}

\begin{theorem}\label{thm:qa-2}
Если в дополнение к условиям теоремы~\ref{thm:qa-1} функция $\varphi(x)$
\textbf{сильно выпукла}, то есть
\[
  \exists\, m, M > 0,\ m \le M : \qquad
  m \norm{y}^2 \le y\T H(x)\, y \le M \norm{y}^2 ,
\]
то $\{x_k\} \to x_*$ со \textbf{сверхлинейной} скоростью.
\end{theorem}

\begin{idea}[title={Доказательство (на лекции не приводилось)}]
\begin{proof}
\textbf{Шаг 1: $x_k \to x_*$.} Из сильной выпуклости
$\tfrac{m}{2}\norm{x_k - x_*}^2 \le \varphi(x_k) - \varphi(x_*) \to 0$ по
теореме~\ref{thm:qa-1}, значит $x_k \to x_*$ (точка $x_*$ при сильной
выпуклости единственна).

\medskip
\textbf{Шаг 2: единичный шаг в конце концов принимается.} По
лемме~\ref{lem:qa-basic}\,$3^\circ$
$\norm{s_k}^2 \le \tfrac{2}{m}\abs{\eta_k} \to 0$. По формуле Тейлора
\[
  \varphi(x_k + s_k) - \varphi(x_k)
  = \nabla\T\varphi(x_k)s_k + \tfrac12 s_k\T H(x_k) s_k + o\big(\norm{s_k}^2\big)
  = \eta_k + o\big(\norm{s_k}^2\big) ,
\]
так что условие~\eqref{eq:qa-armijo} при $\alpha = 1$ равносильно
\[
  (1 - \eps)\,\abs{\eta_k} \ \ge\ o\big(\norm{s_k}^2\big) .
\]
Но $\abs{\eta_k} \ge \tfrac{m}{2}\norm{s_k}^2$, то есть левая часть имеет
порядок $\norm{s_k}^2$, а правая — строго меньший. Значит начиная с некоторого
номера $\alpha_k = \alpha_0 = 1$, и
\[
  x_{k+1} = x_k + s_k = y_k .
\]

\medskip
\textbf{Шаг 3: оценка расстояния.} Обозначим
$e_k = x_k - x_*$, $d_k = y_k - x_*$. Точка $y_k$ минимизирует
$q_k(x) = \varphi(x_k) + \psi_k(x)$ на выпуклом $S$, поэтому
\[
  \nabla\T q_k(y_k)\,(x - y_k) \ge 0 \quad \forall x \in S ,
  \qquad
  \nabla q_k(y) = \nabla\varphi(x_k) + H(x_k)(y - x_k) ,
\]
а точка $x_*$ минимизирует $\varphi$ на $S$:
$\nabla\T\varphi(x_*)(x - x_*) \ge 0$ $\forall x \in S$. Подставим в первое
неравенство $x = x_*$, во второе $x = y_k$ и сложим:
\[
  \big[\nabla\varphi(x_k) - \nabla\varphi(x_*) + H(x_k)(y_k - x_k)\big]\T
  (x_* - y_k) \ \ge\ 0 .
\]
По формуле Ньютона--Лейбница
$\nabla\varphi(x_k) - \nabla\varphi(x_*) = G_k e_k$, где
$G_k = \int_0^1 H(x_* + t e_k)\,dt$. Так как $y_k - x_k = d_k - e_k$,
последнее неравенство переписывается как
\[
  \big[H(x_k)\, d_k + \big(G_k - H(x_k)\big) e_k\big]\T d_k \ \le\ 0 ,
\]
откуда
\[
  m \norm{d_k}^2 \ \le\ d_k\T H(x_k)\, d_k
  \ \le\ \norm{G_k - H(x_k)} \cdot \norm{e_k} \cdot \norm{d_k} ,
\]
то есть
\begin{equation}\label{eq:qa-rate}
  \norm{y_k - x_*} \ \le\ \frac{1}{m}\,\norm{G_k - H(x_k)}\cdot\norm{x_k - x_*} .
\end{equation}

\medskip
\textbf{Шаг 4.} Гессиан непрерывен, а $x_k \to x_*$, поэтому
\begin{align*}
  \norm{G_k - H(x_k)}
  &= \Big\| \int_0^1 \big[H(x_* + t e_k) - H(x_k)\big] dt \Big\| \\
  &\le \sup_{t \in [0,1]} \norm{H(x_* + t e_k) - H(x_k)}
   \ \to\ 0 , \quad k \to \infty .
\end{align*}
С учётом шага~2 ($x_{k+1} = y_k$) и~\eqref{eq:qa-rate}
\[
  \norm{x_{k+1} - x_*} = o\big(\norm{x_k - x_*}\big) ,
\]
то есть сходимость сверхлинейная.
\end{proof}
\end{idea}

\begin{theorem}\label{thm:qa-3}
Если дополнительно к условиям теоремы~\ref{thm:qa-2} гессиан $H(\cdot)$
удовлетворяет условию Липшица,
\[
  \norm{H(x) - H(y)} \le R\,\norm{x - y} \qquad \forall x, y ,
\]
то скорость сходимости \textbf{квадратичная}:
\[
  \norm{x_{k+1} - x_*} \ \le\ \frac{R}{2m}\,\norm{x_k - x_*}^2 .
\]
\end{theorem}

\begin{idea}[title={Доказательство (на лекции не приводилось)}]
\begin{proof}
Повторяем шаги~1--3 доказательства теоремы~\ref{thm:qa-2} вплоть
до оценки~\eqref{eq:qa-rate}; остаётся уточнить $\norm{G_k - H(x_k)}$.
Точка $x_* + t e_k$ отстоит от $x_k = x_* + e_k$ на $(1-t)\norm{e_k}$, поэтому
по условию Липшица
\[
  \norm{G_k - H(x_k)}
  \ \le\ \int_0^1 \norm{H(x_* + t e_k) - H(x_k)}\,dt
  \ \le\ \int_0^1 R\,(1 - t)\,\norm{e_k}\,dt
  = \frac{R}{2}\,\norm{e_k} .
\]
Подставляя в~\eqref{eq:qa-rate} и пользуясь $x_{k+1} = y_k$,
\[
  \norm{x_{k+1} - x_*} \ \le\ \frac{R}{2m}\,\norm{x_k - x_*}^2 . \qedhere
\]
\end{proof}
\end{idea}

\begin{idea}
Три теоремы образуют лестницу, типичную для методов второго порядка:
\begin{itemize}[nosep,left=0pt]
  \item выпуклость $+$ дробление шага дают \emph{глобальную} сходимость по
        значению функции, но без оценки скорости;
  \item сильная выпуклость добавляет единственность $x_*$, сходимость самих
        точек и \emph{сверхлинейную} скорость — за счёт того, что вблизи
        минимума единичный шаг перестаёт отвергаться;
  \item липшицевость гессиана превращает сверхлинейную оценку в
        \emph{квадратичную}, в точности как в теореме~\ref{thm:newton-rate}
        для безусловной задачи.
\end{itemize}
Содержательно вся разница с методом Ньютона в том, что минимум квадратичной
модели берётся не по всему $\R^n$, а по $S$; ключевое
неравенство~\eqref{eq:qa-rate} получается не из явной формулы
$y_k = x_k - H^{-1}\nabla\varphi$, а из вариационного неравенства — условия
минимума на выпуклом множестве.
\end{idea}
